%==========================================================================
% Artigo (versão final) no template UNESC
% Baseado em: Exemplo de artigo cientifico UNESC (v1.0 - 25/06/2025)
%==========================================================================

%==========================================================================
% Pacotes utilizados pela UNESC - não delete as linhas abaixo
%==========================================================================
\documentclass{UNESC}
\usepackage[T1]{fontenc}
\usepackage[utf8]{inputenc}%\usepackage[latin1]{inputenc}
\renewcommand{\rmdefault}{phv} % Arial
\renewcommand{\sfdefault}{phv} % Arial
\usepackage[brazilian]{babel}
\usepackage{graphicx}
%\usepackage{subfigure}
\usepackage{hyperref}
\usepackage{quoting}
\usepackage[labelsep=none]{caption}
\DeclareCaptionLabelFormat{figuracomtraco}{#1~#2 -}
\captionsetup[figure]{labelformat=figuracomtraco, labelsep=space}
\captionsetup[table]{labelformat=simple, labelsep=space, textformat=simple}
\DeclareCaptionLabelFormat{tabelacomtraco}{#1~#2 -}
\captionsetup[table]{labelformat=tabelacomtraco}
\captionsetup{labelsep=none}
%\renewcommand{\thefigure}{\arabic{figure} - }
\usepackage{enumitem}
\usepackage{float}
\usepackage{comment}
\usepackage[hang,flushmargin]{footmisc}
\usepackage{array}

\hypersetup{
  setpagesize  = false,
  colorlinks   = true,    % Colours links instead of ugly boxes
  urlcolor     = blue,    % Colour for external hyperlinks
  linkcolor    = black,   % Colour of internal links
  citecolor    = black    % Colour of citations
}
\urlstyle{same}
\usepackage{setspace}

\usepackage[alf,recuo=0.0cm,abnt-emphasize=bf]{abntex2cite} 

%\captionsetup[figure]{justification=raggedright}


%==========================================================================
% Dados do artigo: título, autores e filiação
%==========================================================================
\title{DETECÇÃO DE REGRESSÕES VISUAIS EM INTEGRAÇÃO CONTÍNUA: IMPLEMENTAÇÃO E ANÁLISE DE TÉCNICAS PARA DETECTAR FALHAS DE INTERFACE}
\author
{
   Douglas Martignago Rosso\footnote{Curso de Ciência da Computação, Universidade do Extremo Sul Catarinense (Unesc), Criciúma - Santa Catarina - Brasil. douglasrosso10@gmail.com},
   Ana Cláudia Garcia Barbosa\footnote{Orientadora. Curso de Ciência da Computação, Universidade do Extremo Sul Catarinense (Unesc), Criciúma - Santa Catarina - Brasil. agb@unesc.net}
}
%==========================================================================
% Início do texto do artigo
%==========================================================================
\begin{document}

\onehalfspacing

\maketitle

\noindent
\textbf{Resumo:} Este trabalho avalia técnicas de regressão visual em integração contínua, comparando diferença \emph{pixel a pixel}, métrica perceptual SSIM e abordagem por regiões com uso de máscaras. A verificação visual foi automatizada em \emph{pipelines} de CI/CD, com \emph{baseline} sob controle de versão e geração de artefatos rastreáveis (relatório HTML e interface de revisão). O estudo discute capacidade de detecção, ocorrência de alertas indevidos e esforço de triagem, considerando variações de renderização, alterações de \emph{layout}, conteúdo dinâmico e responsividade.

\vspace{0.5em}
\noindent
\textbf{Palavras-chave:} regressão visual; integração contínua; testes automatizados; SSIM; baseline.

\vspace{1.0em}
\noindent
\textbf{ABSTRACT:} This work evaluates visual regression techniques in continuous integration by comparing pixel-by-pixel diff, the SSIM perceptual metric, and a region-based approach with masking. Visual checks were automated in CI/CD pipelines, with baselines under version control and traceable artifacts (HTML report and a review UI). The study discusses detection capability, noisy alerts, and triage effort under rendering variability, layout changes, dynamic content, and responsiveness.

\vspace{0.5em}
\noindent
\textbf{Keywords:} visual regression; continuous integration; automated testing; SSIM; baseline.


\section{INTRODUÇÃO}

Em ambientes corporativos, a necessidade de garantir qualidade em ciclos curtos de desenvolvimento incentiva a adoção de testes automatizados em áreas de maior risco e impacto. Conforme destacam \citeonline{costa2022}, essa prática aumenta eficiência, confiabilidade e velocidade de validação, favorecendo entregas contínuas e estáveis.

As interfaces evoluem com frequência. Mudanças de posicionamento de elementos, ajustes tipográficos ou inclusão de componentes podem alterar a apresentação visual sem modificar a lógica. Quando a versão atual diverge da referência aprovada, caracteriza-se uma regressão visual, e a verificação compara a captura recente com a imagem de referência para decidir se a alteração afeta padrões de qualidade, como descrevem \citeonline{tanno2020}.

A comparação \emph{pixel a pixel} é sensível a variações mínimas, porém, como observa \citeonline{tanno2020}, tende a produzir alertas diante de deslocamentos globais e flutuações de renderização. Para mitigar esse ruído, são úteis a segmentação por regiões, a correspondência entre versões e o uso de máscaras em elementos naturalmente dinâmicos. Além disso, métricas perceptuais como a SSIM aproximam a comparação da percepção humana ao considerar luminância, contraste e estrutura \cite{wang2004,wang2003msssim}.

Na prática, o teste visual integra o fluxo de Integração e Entrega Contínuas, com geração de artefatos e vinculação dos resultados a \emph{commits}, evitando que diferenças indesejadas avancem no processo de construção. Revisões e estudos observacionais, como os de \citeonline{shahin2017} e \citeonline{hilton2016ci}, associam CI/CD a ciclos mais frequentes e retorno de informação mais rápido \cite{regoneto2021,silva2023ic}, o que reforça o papel de pontos automatizados de checagem de qualidade em pipelines \cite{maia2024}.

Esse tipo de verificação é compatível com metodologias ágeis, como Scrum e Kanban, que valorizam fluxo contínuo, entregas incrementais e \emph{feedback} frequente \cite{bessa2018}. Relatos nacionais, como o de \citeonline{costa2023}, indicam ganhos de qualidade e produtividade com a transformação ágil. Assim, este artigo avalia três técnicas de regressão visual: comparação \emph{pixel a pixel}, comparação perceptual por SSIM e abordagem por regiões, considerando detecção de defeitos, ruído de alertas indevidos, custo operacional e adequação a mudanças de estilo, à distribuição de elementos, ao conteúdo dinâmico e à responsividade. A avaliação foi automatizada em \emph{pipelines} de CI/CD: a cada execução, tarefas geraram e compararam capturas, aplicaram as métricas e máscaras definidas, consolidaram resultados e publicaram relatórios. As referências permaneceram sob controle de versão, garantindo auditoria e rastreabilidade ao longo da evolução do código.

\vspace{0.5em}
\noindent
\textbf{Objetivo geral:} Avaliar técnicas de regressão visual em equipes ágeis quanto à detecção de falhas, à ocorrência de falsos positivos e ao esforço de triagem, com imagens de referência sob controle de versão na CI/CD.

\vspace{0.5em}
\noindent
\textbf{Objetivos específicos:}
\begin{itemize}[leftmargin=2em]
    \item Identificar cenários e tipos de mudança na interface com maior risco de regressão visual.
    \item Implementar um fluxo automatizado de verificação em CI/CD para geração de artefatos rastreáveis.
    \item Mensurar desempenho das técnicas em detecção, falsos positivos e esforço de triagem.
    \item Analisar robustez diante de variações de renderização, \emph{layout} e conteúdo dinâmico.
\end{itemize}

\vspace{0.5em}
\noindent
\textbf{Justificativa:} A automação de testes amplia cobertura, reduz esforço manual e mitiga erros quando há registro de evidências e análise de \emph{logs} que sustentam a governança da qualidade \cite{alberto2020,banerjee2013}. Em organizações com integração e entrega contínuas, decisões de liberação dependem de \emph{feedback} rápido e rastreável, o que reforça o papel dos testes automatizados no fluxo \cite{shahin2017}. Do ponto de vista prático, combinar automação, evidências reproduzíveis e visibilidade do processo favorece decisões técnicas mais informadas.

Em sistemas de maior porte, manter testes de regressão exclusivamente manuais é caro e pouco escalável. Estudos de caso, como o de \citeonline{costa2022}, recomendam priorizar a automação por risco, criticidade e esforço de manutenção. No recorte visual, imagens de diferenças \emph{pixel a pixel} aumentam o número de falsos positivos diante de variações sutis e deslocamentos, enquanto medidas perceptuais como a SSIM atenuam esse efeito \cite{wang2004,ding2022}, e técnicas \emph{region-based} segmentam e casam regiões para salientar diferenças essenciais. Em áreas dinâmicas, máscaras reduzem alertas desnecessários sem ocultar defeitos relevantes \cite{tanno2020}. Em cenários reais, a calibragem de limiares (valor numérico de tolerância), máscaras e segmentação tende a produzir relatórios úteis para a triagem.

Segundo \citeonline{shahin2017,hilton2016ci,regoneto2021,lima2021,silva2023ic}, integrar a verificação visual e outros testes automatizados ao \emph{pipeline} favorece rastreabilidade, pois imagens de referência e relatórios ficam vinculados a \emph{commits}, \emph{branches} e solicitações de mudança \cite{maia2024}. Essa integração exige ambiente determinístico (sistema operacional, fontes e \emph{renderers}), critérios explícitos de aceitação e políticas de atualização de \emph{baseline} (imagem de referência aprovada que serve como padrão de comparação).

No plano de implantação, recomenda-se, conforme \citeonline{heinonen2020}, estabilizar dados de teste, definir sementes (valores de inicialização) e padronizar o ambiente, a fim de evitar instabilidades e elevar a previsibilidade do \emph{pipeline}. Essa abordagem é coerente com Scrum e Kanban, que priorizam fluxo contínuo, entregas incrementais e \emph{feedback} frequente \cite{bessa2018}. Nesse contexto, testes visuais automatizados funcionam como \emph{gate} (ponto de checagem/critério de aprovação) de qualidade entre versões, sem comprometer a cadência, quando bem calibrados.

O objetivo desta análise foi produzir evidências comparativas para escolher e integrar ao CI/CD a opção técnica que melhor equilibre capacidade de detecção de defeitos, ruído de falsos positivos, esforço de triagem e robustez frente a variações de renderização, \emph{layout} e conteúdo dinâmico. Para isso, foram confrontadas três estratégias representativas, a saber, comparação \emph{pixel a pixel}, comparação perceptual baseada em SSIM e correspondência por regiões com uso de máscaras, com coleta padronizada de métricas e geração de artefatos rastreáveis. A partir dos achados, consolidaram-se diretrizes práticas para operação, incluindo recomendações sobre máscaras, segmentação, manutenção de \emph{baseline} e limiares, além de pontos de atenção no desenho de captura para reduzir instabilidades.

\vspace{0.5em}
\noindent
\textbf{Estrutura do trabalho:} Este artigo está organizado em cinco seções. Na Introdução, apresenta-se o problema de regressão visual, a motivação para automação e a integração com CI/CD, além dos objetivos e da justificativa. Em seguida, as bases conceituais são desenvolvidas na Fundamentação Teórica. A seção de Trabalhos Correlatos compara soluções e relatos da literatura, destacando lacunas. Em Materiais e Métodos, descrevem-se o desenho experimental, as técnicas analisadas, as métricas e a configuração do fluxo, incluindo diferentes larguras de visualização. Por fim, em Resultados e Discussão e na Conclusão, analisam-se os achados e sintetizam-se recomendações.


\section{FUNDAMENTAÇÃO TEÓRICA}

Esta seção reúne os conceitos que sustentam o estudo: princípios de testes e automação; definição de regressão visual e comparação entre três famílias de técnicas (diferença \emph{pixel a pixel}, métricas perceptuais SSIM/MS-SSIM e abordagens por regiões), com tratamento de conteúdo dinâmico, manutenção de \emph{baseline} e necessidade de determinismo; integração em \emph{pipelines} de CI/CD; e a relação com métodos ágeis.

\subsection{Testes de software e automação}

De acordo com \citeonline{alberto2020} e \citeonline{banerjee2013}, testes de software contribuem para qualidade, confiabilidade e manutenção. A automação aumenta cobertura e repetibilidade e apoia a governança do processo. Em escala maior, prioriza-se automatizar onde risco e impacto são mais altos, viabilizando verificações frequentes em ciclos curtos \cite{costa2022}.

\subsection{Regressão visual}

Regressão visual é a diferença indesejada na apresentação da interface entre versões, identificada por comparação de imagens de tela \cite{tanno2020,figueiredo2018,bordignon2019,tanoue2018vrt}. O objetivo é indicar se as mudanças afetam padrões de qualidade previamente definidos, conforme discutem \citeonline{tanno2020}.

\subsubsection{Comparação pixel a pixel}
A diferença \emph{pixel a pixel} marca qualquer variação entre as imagens e, por isso, revela alterações sutis. Contudo, conforme destacam \citeonline{tanno2020,figueiredo2018,tanoue2018vrt}, flutuações de renderização, ajustes de fonte ou diferenças de \emph{layout} podem elevar o número de alertas e aumentar o esforço de triagem.

\subsubsection{Métricas perceptuais: SSIM e MS-SSIM}
Métricas perceptuais aproximam a avaliação da percepção humana. A SSIM, proposta por \citeonline{wang2004}, compara luminância, contraste e estrutura e costuma ser mais tolerante a variações pequenas que a diferença \emph{pixel a pixel}. A MS-SSIM estende a análise a múltiplas escalas, o que favorece robustez a diferentes resoluções e níveis de detalhe \cite{wang2003msssim}. Na prática, limiares de aceitação e máscaras são combinados para equilibrar sensibilidade e ruído.

\subsubsection{Abordagens por regiões}
Abordagens por regiões segmentam a tela, realizam correspondência entre partes equivalentes e enfatizam apenas diferenças essenciais. Essa estratégia, segundo \citeonline{tanno2020}, reduz o impacto de deslocamentos e mudanças de renderização, tornando as imagens de diferenças mais informativas para a triagem.

\subsubsection{Conteúdo dinâmico, baseline e determinismo}
Elementos naturalmente variáveis (datas, contadores, banners, dados externos) produzem diferenças recorrentes se não tratados. Conforme indicam \citeonline{tanno2020} e \citeonline{tanoue2018vrt}, recomenda-se aplicar máscaras apenas em regiões inevitavelmente variáveis e adotar política explícita de atualização de \emph{baseline} para manter rastreabilidade e evitar normalização de defeitos. Como a comparação ocorre entre execuções distintas, o determinismo operacional é essencial: versões fixas de sistema, navegador e controladores, resolução e densidade constantes, fontes empacotadas, idioma e fuso horário controlados, animações desativadas e simulação de dados quando necessário \cite{heinonen2020}.

\subsection{Integração a pipelines de CI/CD}

Integrar a verificação visual ao fluxo de CI/CD aumenta a rastreabilidade e acelera o retorno de informação: imagens de referência, diferenças e relatórios vinculam-se a \emph{commits}, \emph{branches} e solicitações de mudança. Em sua revisão, \citeonline{shahin2017} e \citeonline{hilton2016ci} associam integrações frequentes a ciclos mais curtos e \emph{feedback} rápido, o que favorece a adoção de pontos automatizados de qualidade \cite{maia2024}.

\subsection{Métodos ágeis e o papel do feedback}

Em Scrum e Kanban, fluxo contínuo, entregas incrementais e retorno frequente de informação sustentam verificações visuais leves entre versões \cite{bessa2018}. Estudos nacionais, como o de \citeonline{costa2023}, relacionam a transformação ágil a ganhos de qualidade e produtividade, o que reforça a aderência prática dessa abordagem.

\subsection{Síntese e lacuna}

Há espaço para comparar, sob condições controladas, três classes de técnicas, diferença \emph{pixel a pixel}, SSIM/MS-SSIM e abordagens por regiões, considerando detecção, alertas indevidos e esforço de triagem, com integração efetiva ao CI/CD e artefatos rastreáveis \cite{tanno2020,wang2004,wang2003msssim,heinonen2020}. Este trabalho aborda essa lacuna com um experimento em diferentes larguras de visualização e com coleta padronizada de métricas e evidências.


\section{TRABALHOS CORRELATOS}

A literatura nacional e internacional oferece bases técnicas para testes visuais e sua operação em fluxos de desenvolvimento. A seguir, sintetizam-se alguns trabalhos que dialogam com este estudo quanto à comparação por imagem, uso de métricas perceptuais e prática em Integração/Entrega Contínuas.

Em mapeamento sistemático sobre testes de interfaces gráficas, \citeonline{banerjee2013} discutem abordagens, ferramentas e desafios para automação de GUI, destacando a importância de aumentar a cobertura e reduzir o esforço manual em cenários complexos. Em um estudo de caso nacional, \citeonline{alberto2020} relatam a implantação de automação de testes em um contexto organizacional real, descrevendo ganhos de rastreabilidade, padronização do processo e redução de retrabalho, o que reforça a relevância prática de estratégias automatizadas de verificação.

\citeonline{wang2004} introduzem a SSIM como métrica perceptual que considera luminância, contraste e estrutura, aproximando a avaliação da percepção humana e reduzindo falsos positivos em relação à diferença \emph{pixel a pixel}. Posteriormente, \citeonline{ding2022} propõem uma formulação que unifica estrutura e textura na avaliação de qualidade de imagem, enquanto \citeonline{zhang2018lpips} exploram o uso de \emph{deep features} como métricas perceptuais mais alinhadas ao julgamento humano. Esses trabalhos oferecem base conceitual para empregar medidas perceptuais em testes de regressão visual.

Em revisão extensa, \citeonline{shahin2017} sintetizam abordagens, ferramentas, desafios e práticas de Integração/Entrega Contínuas, relacionando integrações frequentes a ciclos mais curtos e \emph{feedback} acelerado. Complementarmente, \citeonline{hilton2016ci} analisam o uso de integração contínua em projetos de código aberto, evidenciando benefícios e custos associados à automação. No contexto nacional, estudos aplicados como os de \citeonline{maia2024} descrevem a construção de pipelines de CI/CD com testes automatizados em sistemas reais, reforçando a importância de ligar artefatos de teste a \emph{commits}, padronizar o ambiente de execução e utilizar a esteira como ponto de checagem de qualidade.

Em aplicações Android, \citeonline{figueiredo2018} descrevem a utilização de comparação de imagens em um processo de testes regressivos visuais, combinando automação e testes exploratórios em um ambiente ágil. De forma complementar, \citeonline{bordignon2019} realizam um estudo de caso com ferramentas de \emph{Visual GUI Testing} para verificar requisitos não funcionais em interfaces gráficas, destacando desafios práticos semelhantes aos relatados na literatura internacional. Ainda em âmbitos nacionais, \citeonline{lima2021} sintetizam a adoção de técnicas de teste de regressão automatizado em projetos Android, enquanto \citeonline{quirino2025} analisam empiricamente os \emph{plugins} de teste mais utilizados em Pytest, Jest e Cypress, incluindo extensões de regressão visual —, e \citeonline{regoneto2021,silva2023ic} investigam a relação entre integração contínua e qualidade de software em projetos industriais.


\section{MATERIAIS E MÉTODOS}

Esta pesquisa é aplicada, pois compara técnicas de detecção de regressão visual e as integra, de forma prática, a fluxos de Integração e Entrega Contínuas (CI/CD). Em termos de objetivos, é exploratória e descritiva: exploratória por ampliar o entendimento sobre abordagens e parâmetros de calibração; descritiva por relatar, de modo organizado, os procedimentos adotados e os resultados obtidos. Quanto aos procedimentos, utiliza experimentação controlada com análise quantitativa, baseada em medidas numéricas, e síntese descritiva. Como base conceitual, adotam-se métricas perceptuais de imagem, notadamente a SSIM, proposta por \citeonline{wang2004}, e práticas de operação em pipelines de software discutidas por \citeonline{shahin2017} e \citeonline{heinonen2020}.

\subsection{Objetivo metodológico}
O objetivo é identificar, em condições controladas, qual técnica de verificação visual oferece melhor equilíbrio entre quatro aspectos: capacidade de detectar alterações realmente relevantes; incidência de alertas indevidos; esforço de triagem, medido como o tempo gasto para analisar cada alerta; e robustez diante de variações comuns de renderização, distribuição dos elementos na tela, conteúdo dinâmico e diferentes larguras de visualização. A motivação técnica para incluir uma métrica perceptual decorre do fato de que, conforme \citeonline{wang2004}, a SSIM considera luminância, contraste e estrutura, aproximando a comparação da percepção humana.

\subsection{Ambiente experimental e determinismo}
Os testes foram executados em ambiente padronizado para garantir resultados reproduzíveis. Utilizou-se Node.js (v20) e Playwright (Chromium \emph{headless}) para captura de telas, com densidade de pixels fixa e \emph{locale} \texttt{pt-BR}/fuso \texttt{America/Sao\_Paulo}. Durante a captura, animações e transições foram desativadas e, para reduzir instabilidades, o relógio e a função \texttt{Math.random} foram congelados por \emph{script} injetado antes do carregamento da página. As capturas foram feitas com \emph{screenshot} \emph{full-page}, e as imagens atuais foram salvas em \texttt{results/current/}. As imagens de referência (\emph{baseline}) permaneceram em \texttt{baselines/} sob controle de versão, permitindo auditoria e rastreabilidade. Esses cuidados de determinismo operacional são recomendados para reduzir ruído em verificações automatizadas no \emph{pipeline} \citeonline{heinonen2020}.

\subsection{Delineamento experimental}
O delineamento especifica os elementos do experimento e como eles se combinam:
\begin{enumerate}[label=\alph*),leftmargin=1.5em]
  \item Cenário avaliado. O cenário central foi uma aplicação \emph{single-page} em React (tela \textit{dashboard}), capturada de forma determinística (congelamento de relógio e \texttt{Math.random}) para reduzir variações entre execuções. Apesar disso, registrou-se também um caso relevante de instabilidade: diferenças de altura em capturas \emph{full-page} (mesma largura, altura levemente diferente), que afetam diretamente a comparabilidade por imagem.
  \item Técnicas. Três famílias representativas foram comparadas: (i) comparação \emph{pixel a pixel} com tolerância (\textit{pixelmatch}); (ii) comparação perceptual baseada em SSIM, conforme \citeonline{wang2004}, com implementação em blocos (\emph{blockSize}); e (iii) abordagem por regiões com segmentação em grade e suporte a máscaras para áreas dinâmicas \citeonline{tanno2020}.
  \item Larguras de visualização. Foram utilizados três \emph{viewports} representativos: 360\,px (\textit{mobile}), 768\,px (\textit{tablet}) e 1366\,px (\textit{desktop}). Cada página foi capturada em todas as larguras.
  \item Máscaras e política de aceitação. A técnica por regiões foi configurada para grade 4\,$\times$\,6, com suporte a mascaramento por célula. Neste recorte, as máscaras permaneceram desativadas (lista vazia) para evidenciar o comportamento de ruído/alertas. A atualização de \emph{baseline} foi tratada como processo explícito e rastreável no fluxo de CI/CD.
\end{enumerate}

\subsection{Implementação do fluxo (PixelGuard)}
Para viabilizar a avaliação e a operação em \emph{pipeline}, foi implementada uma ferramenta de regressão visual (PixelGuard) composta por três etapas automatizadas: captura, comparação e publicação de evidências.

Na etapa de captura (\texttt{tests/capture.js}), a aplicação React é servida localmente via Vite e o Playwright realiza capturas \emph{full-page} em três \emph{viewports}. Para reduzir instabilidades, o script injeta um \emph{init script} que congela a data e um gerador pseudoaleatório determinístico, além de desativar animações e transições.

Na etapa de comparação (\texttt{tests/compare.js}), as imagens atuais (\texttt{results/current/}) são confrontadas com as baselines (\texttt{baselines/}). Cada imagem é processada por três comparadores: (i) pixel a pixel (\texttt{pixelmatch}) com tolerância \texttt{0.1} e limite de \texttt{0.05\%} de diferença; (ii) SSIM com limiar mínimo \texttt{0.999} e \emph{blockSize} \texttt{8}; e (iii) regiões com grade 4\,$\times$\,6 e limite de \texttt{1.0\%} por célula. O resultado consolidado é salvo em \texttt{results/results.json}.

Na etapa de evidências e triagem, um relatório HTML (\texttt{results/report.html}) é gerado com resumo e imagens (baseline, atual e \emph{diffs}). Adicionalmente, uma interface de revisão (\emph{Review UI}) permite inspecionar as diferenças e tomar a decisão de aprovação/rejeição, tornando o processo rastreável e operacionalmente viável.

\subsection{Métricas e coleta}
Foram coletados indicadores gerados automaticamente pelas técnicas e pelo próprio fluxo de execução. Para a comparação \emph{pixel a pixel}, registrou-se o percentual de pixels diferentes (\texttt{diffPercent}) e a contagem de pixels divergentes. Para SSIM, registrou-se o \emph{score} de similaridade estrutural (\texttt{score}) e a proporção de blocos com baixa similaridade (\texttt{diffPercent}). Para a abordagem por regiões, registraram-se a quantidade de células reprovadas e o total de células (\texttt{failedRegions/totalRegions}), além do percentual de diferença por célula. Também foram registrados erros operacionais (p.ex., dimensões diferentes entre \emph{baseline} e captura atual) e o tempo de execução por técnica (\texttt{executionMs}).

Métricas clássicas como taxa de detecção, taxa de alertas indevidos e tempo de triagem cronometrado dependem de um conjunto rotulado de alterações (mudanças injetadas vs. não relevantes) e de protocolo de triagem com coleta de tempo. Como o foco desta execução foi validar o fluxo ponta a ponta e obter evidências rastreáveis por técnica, esses indicadores foram tratados como diretrizes de avaliação e como recomendações para expansão do experimento.

\subsection{Procedimentos de análise}
Os resultados foram comparados entre técnicas e entre as diferentes larguras de visualização (\textit{mobile}, \textit{tablet} e \textit{desktop}). A análise discutiu o comportamento relativo de sensibilidade e a interpretabilidade das evidências (imagens de \emph{diff} e contagens/\emph{scores}), conciliando detecção de mudanças relevantes com redução de ruído operacional. Para interpretar diferenças entre as abordagens, considerou-se a literatura sobre comparação perceptual, como a de \citeonline{wang2004}, e sobre operação em fluxos contínuos de integração e entrega, a exemplo de \citeonline{shahin2017}.

\subsection{Validade e limitações}
Foram documentadas ameaças à validade interna, como viés na seleção de telas e sensibilidade dos limiares de aceitação, e à validade externa, como a possibilidade de generalização para outros tipos de aplicações e ambientes. Também foram registradas dependências do ambiente (fontes, \emph{renderers} e resolução) e as medidas adotadas para impedir que variações não controladas influenciassem os resultados. Conforme destacam \citeonline{tanno2020} e \citeonline{tanoue2018vrt}, o uso de máscaras reduz ruído, mas sua calibragem inadequada pode ocultar defeitos; por isso, quando utilizadas, devem ser restritas a regiões inevitavelmente variáveis e revisadas a cada atualização de \emph{baseline}.


\section{RESULTADOS E DISCUSSÃO}

Esta seção apresenta resultados obtidos a partir da execução automatizada do fluxo de regressão visual, com geração de artefatos em \texttt{results/}. Foram comparadas três capturas da tela \textit{dashboard}, uma por \emph{viewport} (\textit{mobile}, \textit{tablet} e \textit{desktop}), totalizando três comparações consolidadas em \texttt{results/results.json}.

\subsection{Resumo numérico das comparações}
A Tabela \ref{tab:resumo-resultados} sintetiza os indicadores produzidos por cada técnica. Observa-se que diferenças de dimensão (largura/altura) resultam em falha imediata, pois inviabilizam a correspondência direta entre pixels e blocos, comportamento que afeta igualmente as três técnicas.

\begin{table}[H]
\centering
\caption{Resumo das comparações (extraído de \texttt{results/results.json})}
\label{tab:resumo-resultados}
\small
\setlength{\tabcolsep}{6pt}
\renewcommand{\arraystretch}{1.2}
\begin{tabular}{|l|c|c|c|}
\hline
\textbf{Imagem} & \textbf{Pixel} & \textbf{SSIM} & \textbf{Regiões} \\
\hline
dashboard-desktop-1366w & FAIL (1.701\%) & FAIL (0.9457) & FAIL (17/24) \\
dashboard-mobile-360w & FAIL (dimensões) & FAIL (dimensões) & FAIL (dimensões) \\
dashboard-tablet-768w & FAIL (dimensões) & FAIL (dimensões) & FAIL (dimensões) \\
\hline
\end{tabular}
\end{table}

No caso \textit{desktop}, a comparação \emph{pixel a pixel} registrou 1.701\% de pixels diferentes, enquanto a SSIM obteve \emph{score} 0.9457 e a abordagem por regiões identificou 17 regiões com falha em um total de 24. Esse resultado ilustra que, diante de mudanças visuais mais amplas (por exemplo, alteração de \emph{layout} ou estilo), as três técnicas convergem para reprovação, ainda que com formas distintas de evidenciar o problema: (i) difusão de diferenças em pixels, (ii) degradação do \emph{score} perceptual e (iii) concentração espacial por células.

Para \textit{mobile} e \textit{tablet}, o erro reportado foi \textit{Dimensões diferentes} (mesma largura, altura diferente por poucos pixels). Na prática, isso ocorre quando a captura \emph{full-page} produz alturas diferentes por pequenas variações de conteúdo renderizado, \emph{breakpoints} e/ou arredondamentos de layout. O achado é relevante porque revela uma limitação operacional: antes de discutir \emph{thresholds} e máscaras, o processo precisa garantir um critério consistente de enquadramento (p.ex., capturas por área fixa, recorte por componente/\emph{bounding box} ou normalização de altura) para manter comparabilidade.

\subsection{Evidências visuais e triagem}
O fluxo gerou imagens de diferenças para as três técnicas (pixel, SSIM e regiões), além de um relatório HTML com resumo e evidências. A Figura \ref{fig:report-html} exemplifica o formato do relatório, que consolida o resultado por técnica e apresenta as imagens necessárias para inspeção.

\begin{figure}[H]
\centering
\includegraphics[width=0.95\linewidth]{Figs/09-report-html.png}
\caption{Relatório HTML gerado automaticamente com resumo e evidências visuais}
\label{fig:report-html}
\end{figure}

Além do relatório, a interface de revisão (\emph{Review UI}) foi utilizada para suportar triagem, permitindo selecionar uma captura e alternar entre técnicas de \emph{diff}. A Figura \ref{fig:review-overview} ilustra a tela de revisão, que organiza lista de diferenças, visualização e painel de decisão.

\begin{figure}[H]
\centering
\includegraphics[width=0.95\linewidth]{Figs/01-review-overview.png}
\caption{Interface de revisão (Review UI) para triagem de regressões visuais}
\label{fig:review-overview}
\end{figure}

\subsection{Discussão: sensibilidade, ruído e operação em CI/CD}
Os resultados reforçam aspectos frequentemente relatados na literatura. A comparação \emph{pixel a pixel} é altamente sensível e tende a sinalizar diferenças pequenas, o que pode aumentar falsos positivos quando há instabilidades de renderização \citeonline{tanno2020}. A SSIM, por capturar similaridade estrutural, tende a se comportar de forma mais alinhada ao julgamento humano em variações leves \citeonline{wang2004}, embora ainda dependa de limiares e do modo de agregação (p.ex., por blocos). A abordagem por regiões favorece a triagem ao isolar espacialmente áreas afetadas e permite mascarar apenas regiões inevitavelmente variáveis, reduzindo ruído sem ocultar defeitos relevantes quando bem calibrada \citeonline{tanno2020}.

No plano operacional, a integração em CI/CD foi estruturada em GitHub Actions, com extração de baselines da \textit{main}, execução de captura/comparação e publicação de evidências (artefatos e página de revisão). Um ponto observado na prática é que a estratégia de captura (\emph{full-page}) pode introduzir falhas por divergência de dimensões, o que sugere que a qualidade do \emph{gate} visual depende não apenas da métrica escolhida, mas também do desenho de captura e das garantias de determinismo.


\section{CONCLUSÃO}

Este trabalho avaliou e operacionalizou três técnicas de regressão visual em um fluxo automatizado compatível com CI/CD: diferença \emph{pixel a pixel}, SSIM e abordagem por regiões com suporte a máscaras. A implementação (PixelGuard) demonstrou a viabilidade de manter baselines sob controle de versão, executar capturas determinísticas e publicar evidências rastreáveis (relatório e interface de revisão), fornecendo suporte à tomada de decisão em \emph{pull requests}.

Como principal achado prático, observou-se que a comparabilidade por imagem pode ser comprometida por divergências de dimensão em capturas \emph{full-page}. Esse tipo de instabilidade causa falha imediata em todas as técnicas e, portanto, deve ser tratado como requisito de engenharia do teste (critério de enquadramento e normalização) antes da calibração fina de limiares. No caso em que as dimensões coincidiram (\textit{desktop}), as três técnicas foram capazes de identificar diferenças, com destaque para a abordagem por regiões ao fornecer indicação espacial das áreas afetadas, e para a SSIM ao resumir degradação estrutural em um \emph{score} perceptual.

Como trabalhos futuros, recomenda-se: (i) adotar capturas por componente/\emph{bounding box} ou por altura fixa para reduzir falhas por dimensão; (ii) ampliar o conjunto de páginas e cenários, incluindo conteúdo dinâmico real com máscaras calibradas; (iii) avaliar custos de triagem com medidas de tempo (por técnica e por tipo de mudança); e (iv) comparar variantes perceptuais (p.ex., MS-SSIM) e métricas baseadas em \emph{deep features} em conjunto com políticas de aceitação e atualização de baseline.

\bibliographystyle{abntex2-unesc}
\renewcommand{\refname}{REFERÊNCIAS}
\bibliography{referencias}

\end{document}
